\section{Daten einlesen und Datenbankimport (5123026)}

Das Einlesen der CSV-Dateien erfolgt über ein seperates Batch Skript, welches Python Skripte aufruft. Der Ablauf generiert zuerst die
abhängigen Stations-, Raum- und Buchungsdaten und importiert anschließend alle CSV-Dateien in einer
festen Reihenfolge. Diese Reihenfolge folgt den Fremdschlüsselabhängigkeiten des relationalen
Modells: unabhängige Stammdaten wie \texttt{person} und \texttt{department} werden vor abhängigen
Tabellen wie \texttt{employee}, \texttt{patient}, \texttt{diagnosis} und \texttt{bookings}
geschrieben. Dadurch existieren referenzierte Datensätze bereits, bevor Tabellen mit
Fremdschlüsseln geladen werden.

\subsection{Neuaufbau der Tabellen}

Jeder Import wird mit \texttt{--drop} ausgeführt. Vor dem Einfügen neuer Datensätze wird die
Zieltabelle zurückgesetzt:

\begin{lstlisting}[style=sqlstyle]
TRUNCATE TABLE "public"."<table>" RESTART IDENTITY CASCADE;
\end{lstlisting}

\texttt{TRUNCATE} ist für den Massenimport geeigneter als ein zeilenweises \texttt{DELETE}, weil
die Tabelle als Ganzes geleert wird. \texttt{RESTART IDENTITY} sorgt für reproduzierbare IDs,
\texttt{CASCADE} entfernt abhängige Datensätze mit. So kann der komplette Testdatenbestand
wiederholt sauber neu aufgebaut werden.

\subsection{Schema- und Fremdschlüsselprüfung}

Der Importer liest das aktuelle PostgreSQL-Schema direkt aus den Metadaten. Dafür werden u.\,a. die
Tabellenspalten, Datentypen und benutzerdefinierten Typen abgefragt:

\begin{lstlisting}[style=sqlstyle]
SELECT table_name, column_name, ordinal_position,
       data_type, udt_name, is_nullable
FROM information_schema.columns
WHERE table_schema = %s
ORDER BY table_name, ordinal_position;
\end{lstlisting}

Auf dieser Basis wird das Mapping zwischen CSV-Spalten und Datenbankspalten geprüft. Zusätzlich
werden Enum-Werte, z.\,B. für Buchungsstatus oder Dosierungsfrequenzen, gegen die in PostgreSQL
definierten Werte validiert. Dadurch werden Format- und Typfehler vor dem eigentlichen Insert
sichtbar.

Für die referenzielle Integrität fragt der Importer die Fremdschlüsselbeziehungen des Schemas ab
und prüft anschließend, ob die in der CSV enthaltenen Referenzen bereits in den Elterntabellen
existieren:

\begin{lstlisting}[style=sqlstyle]
SELECT "<foreign_column>"::text
FROM "public"."<foreign_table>"
WHERE "<foreign_column>"::text = ANY(%s);
\end{lstlisting}

Fehlende Referenzen führen zu einem Abbruch vor dem Import. Das ist bei den Krankenhausdaten
besonders wichtig: Buchungen dürfen nur auf vorhandene Patienten und Räume zeigen, Diagnosen nur
auf vorhandene Patienten, Ärzte und Medikationsdaten. Die Vorprüfung verhindert, dass große Batches
erst während des Schreibens an einzelnen Fremdschlüsselverletzungen scheitern.

\subsection{Batch-Inserts und Transaktionen}

Die CSV-Zeilen werden in Paketen von 2.000 Datensätzen importiert. Pro Paket wird ein
Multi-Row-\texttt{INSERT} mit parametrisierten Platzhaltern erzeugt:

\begin{lstlisting}[style=sqlstyle]
INSERT INTO "public"."<table>" ("col_1", "col_2", "col_3")
VALUES (%s, %s, %s), (%s, %s, %s), ...;
\end{lstlisting}

Dadurch werden viele Zeilen mit einer SQL-Anweisung geschrieben, statt für jede CSV-Zeile einen
eigenen Datenbank-Roundtrip auszuführen. Das ist für Tabellen wie \texttt{person}, \texttt{patient},
\texttt{diagnosis} und \texttt{bookings} relevant, da hier mehrere tausend Datensätze verarbeitet
werden.

Jeder Batch läuft in einer Transaktion. Schlägt ein Batch fehl, wird ein Rollback ausgeführt. Zur
Fehleranalyse kann der Importer anschließend auf einzelne Zeilen zurückfallen und Savepoints
verwenden:

\begin{lstlisting}[style=sqlstyle]
SAVEPOINT csv_import_row;
-- einzelner Insert
RELEASE SAVEPOINT csv_import_row;

ROLLBACK TO SAVEPOINT csv_import_row;
\end{lstlisting}

Damit kombiniert der Importprozess zwei Ziele: Im Normalfall werden Massendaten schnell über
Batch-Inserts geschrieben, im Fehlerfall kann trotzdem die konkrete problematische Zeile ermittelt
werden. Für den Use Case der Abgabe ist diese Kombination aus Performance, Transaktionen und
referenzieller Integrität zentral, weil viele voneinander abhängige Tabellen regelmäßig vollständig
neu aufgebaut werden müssen.
